\documentclass{article}
\usepackage{amsmath, amssymb}
\usepackage{amsthm}

\author{Michael Naumov}
\date{}
\setlength{\parindent}{0pt}
\setlength{\parskip}{0.5em}


\title{TAOCP 1.2.1\mbox{-}12}

\begin{document}
\maketitle

\( A := \{ u + v \sqrt{2} | u, v \in \mathbb{Z} \} \)

We modify algorithm \textbf{E} allowing \( m, n, c, d \in A, a, b, a', b', q \in \mathbb{Z} \). We will also modify all invariants from \( > 0 \) to \( \neq 0 \)

All the steps and invariants stay. However, it's a bit unclear how to define calculation of \( q \) deterministically. But it doesn't really matter for the task. We just choose of the possible algorithms to choose \( q \).

Let's prove that the algorithm never ends for \( m = 1, n = \sqrt{2} \)

From the contrary. Assume the algorithm ended. Meaning on some step we got \( r = 0 \)

\( A2 \implies am + bn = d, a'm+b'n = c = qd + r \implies a'm+b'n = qd = q(am+bn) \implies m(a'-qa) = n(qb-b') \)

If we assume \( qb - b' \neq 0 \implies \frac{n}{m} = \frac{a'-qa}{qb-b'} \in \mathbb{Q} \)

\( \frac{n}{m} = \sqrt{2} \notin Q \)

Contradiction.

Then \( qb - b' = 0 \implies a' - qa = 0 \implies b' = qb, a' = qa \)

If \( q = 0 \implies c = r = 0 \), which is impossible, because by \( A2: c \neq 0 \)

Then \( q \neq 0 \)

Let's see what was the result of \( E4 \) that lead to such result.

Let's use with index \( 0 \) to denote variables that were before executing step \( E4 \).

Step \( E4: a' = a_{0}, a = a_{0}' - q_{0}a_{0}, b' = b_{0}, b = b_{0}' - q_{0}b_{0} \)

\( a'=qa \implies a_{0} = q(a_{0}'-q_{0}a_{0}) \implies (1 + qq_{0})a_{0} = q a_{0}' \)

\( a_{0}' = \frac{1+qq_{0}}{q}a_{0}, b_{0}' = \frac{1+qq_{0}}{q}b_{0}, x_{0} := \frac{1+qq_{0}}{q} \)

If \( x_{0} = 0 \implies a_{0}'=b_{0}'=0 \implies a_{0}'m+b_{0}'n=c_{0}=0 \implies \) contradiction \( \implies x_{0} \neq 0 \)

So we see that on previous stage \( a' = xa, b' = xb \) for some \( x \in \mathbb{Q} \setminus \{ 0 \} \)

And repeating this process until the very beginning, we get

\( a'=1, a = 0, a'=xa \implies \) contradiction

All cases are considered. The algorithm never ends.


\end{document}
